% LỜI TỰA --- Quyển IX
\chapter*{Lời Tựa}
\addcontentsline{toc}{chapter}{Lời Tựa}
\markboth{Lời Tựa}{Lời Tựa}

\begin{truyen}{Lời Tựa}{No Draft for Real Life}
\chuhoa{G}{iáo viên bắt đầu bộ sách này với một câu hỏi đơn giản: làm thế nào để mỗi buổi học có thể để lại thứ gì đó đáng nhớ hơn là bài thi?}
\textit{(A teacher began this series with a simple question: how can each class leave behind something more memorable than a test?)}

Ra đời không ai được phép thử trước rồi làm lại~--- đó chính là điều khiến những câu chuyện đầu đời đáng nhớ nhất.
\textit{(No one gets a practice run before real life~--- that is exactly what makes first-life stories the most memorable.)}

Quyển chín dành cho những người đang đứng trước ngưỡng cửa cuộc đời: công việc đầu tiên, đồng tiền đầu tiên tự kiếm, lần đầu bị từ chối, lần đầu sống xa nhà, lần đầu phải tự quyết định tất cả. Không có bản nháp nào cho những khoảnh khắc đó~--- nhưng có những câu chuyện của người đi trước để ta học.
\textit{(Volume nine is for those standing at life's threshold: the first job, the first money earned, the first rejection, the first time living away from home, the first time deciding everything alone. There is no draft for those moments~--- but there are stories from those who went before us to learn from.)}

Bộ sách <<Truyện Tích Cảm Hứng Song Ngữ>> gồm 24 quyển, mỗi quyển một chủ đề riêng. Quyển IX này mang chủ đề: \textbf{công việc đầu tiên, tiền đầu tiên, bị từ chối, sống xa nhà}.
\textit{(The series ``Inspiring Bilingual Tales'' contains 24 volumes, each with its own theme. Volume IX focuses on: \textbf{công việc đầu tiên, tiền đầu tiên, bị từ chối, sống xa nhà}.})
\end{truyen}

\begin{baihoc}
Không có bản nháp nào cho cuộc đời thật~--- nhưng những câu chuyện của người đi trước là bản hướng dẫn tốt nhất.
\textit{(There is no draft for real life~--- but stories from those who came before are the best guides.)}
\end{baihoc}

\begin{ghinhoanh}
\textbf{Short English to remember:}

Life has no rehearsal: show up fully.

First experiences shape who we become.
\end{ghinhoanh}

\begin{tuhocnhanh}
1. Khoảnh khắc nào trong cuộc đời em sắp bước vào mà em lo nhất? \textit{(Which upcoming life moment worries you most?)}

2. Câu chuyện nào trong quyển này giúp em chuẩn bị tâm lý tốt hơn? \textit{(Which story here helped you mentally prepare for what's ahead?)}

3. Em muốn hỏi người đi trước điều gì về cuộc sống sau khi ra trường? \textit{(What would you ask someone ahead of you about life after graduation?)}
\end{tuhocnhanh}

\begin{center}
\textit{\color{truyengold}``Ai đã từng vấp ngã mà đứng dậy được đều có chuyện đáng kể.''}
\end{center}

\begin{flushright}
\textbf{Tôi Nguyễn Văn Sang}\\
\textit{GV THPT Nguyễn Hữu Cảnh - TP HCM}
\end{flushright}

\ngancach
